% =====================================================================
%  NTU-DAVID-RAG — Setup & Deployment Manual
%  Compile:  ./build.sh   (or: pdflatex manual.tex, run twice for ToC)
% =====================================================================
\documentclass[11pt,a4paper]{article}

\usepackage[margin=1in]{geometry}
\usepackage{amssymb}
\usepackage{booktabs}
\usepackage{array}
\usepackage{xcolor}
\usepackage{enumitem}
\usepackage{caption}
\usepackage{float}
\usepackage{listings}
\usepackage[colorlinks=true,linkcolor=teal!60!black,urlcolor=teal!60!black,
            citecolor=teal!60!black]{hyperref}

% ---- section styling (matches the technical report) ----
\usepackage{titlesec}
\definecolor{accent}{RGB}{0,128,128}
\titleformat{\section}{\Large\bfseries\color{accent}}{\thesection}{0.6em}{}
\titleformat{\subsection}{\large\bfseries\color{accent!80!black}}{\thesubsection}{0.6em}{}
\captionsetup{font=small,labelfont={bf,color=accent}}

% ---- shell / code blocks ----
\definecolor{codebg}{RGB}{246,248,248}
\definecolor{codeframe}{RGB}{210,222,222}
\lstdefinestyle{sh}{
  basicstyle=\ttfamily\footnotesize,
  backgroundcolor=\color{codebg},
  frame=single, framesep=5pt, rulecolor=\color{codeframe},
  breaklines=true, breakatwhitespace=false, postbreak=\mbox{\textcolor{accent}{$\hookrightarrow$}\space},
  columns=fullflexible, keepspaces=true, showstringspaces=false,
  xleftmargin=4pt, xrightmargin=4pt, aboveskip=7pt, belowskip=7pt,
  literate={~}{{\textasciitilde}}1
}
\lstset{style=sh}
\newcommand{\code}[1]{\texttt{\small #1}}

\title{\vspace{-1.2cm}\textbf{\color{accent}NTU-DAVID-RAG}\\[2pt]
       \large Setup \& Deployment Manual\\[4pt]
       \normalsize Single-GPU and Two-GPU Docker installations}
\author{}
\date{}

\begin{document}
\maketitle
\vspace{-0.8cm}

\begin{abstract}
\noindent
This manual is a step-by-step guide to installing and running the
NTU-DAVID-RAG service from a clean machine. It covers prerequisites and
GPU/Docker setup for Linux, Windows, and macOS; cloning the repository;
configuring secrets; and two deployment options --- a \textbf{two-GPU}
default (the 7B embedder on one card, the answer LLM on another) and a
\textbf{single-GPU} mode that fits the whole stack on one card. Every step
lists the exact commands.
\end{abstract}

\tableofcontents
\vspace{0.4cm}

% =====================================================================
\section{What you are installing}
% =====================================================================
NTU-DAVID-RAG is a containerized retrieval-augmented-generation service for
economics papers and their Fortran model code. The stack is orchestrated with
Docker Compose and exposes a single web dashboard behind an Nginx proxy. It has
four moving parts:

\begin{itemize}[leftmargin=1.4em,itemsep=2pt]
  \item \textbf{embedding-server} --- a GPU container that loads the text
        embedder and the cross-encoder reranker once and serves them over HTTP.
  \item \textbf{rag-web} --- the Flask/Gunicorn application (ingestion,
        retrieval, generation, auth). It is CPU-only except for MinerU (PDF
        layout/OCR), which uses the GPU transiently during ingestion.
  \item \textbf{answer LLM} --- an Ollama service (e.g.\ \code{gemma4:31b} or
        \code{gemma3:12b}) that writes the grounded answers.
  \item \textbf{auth-db} + \textbf{proxy} --- PostgreSQL for users/sessions and
        an Nginx reverse proxy that publishes the dashboard on one port.
\end{itemize}

\noindent\textbf{Two ways to run it:}
\begin{itemize}[leftmargin=1.4em,itemsep=2pt]
  \item \textbf{Two-GPU (default).} The 7B embedder owns one GPU; the
        \code{gemma4:31b} LLM runs as a separate Ollama stack on another GPU.
        Highest quality. See Section~\ref{sec:twogpu}.
  \item \textbf{Single-GPU.} An override collapses the embedder, reranker,
        MinerU, and an in-stack LLM onto one card, swapping in lighter models.
        See Section~\ref{sec:onegpu}.
\end{itemize}

% =====================================================================
\section{System requirements}\label{sec:req}
% =====================================================================

\subsection{Hardware}
A CUDA-capable NVIDIA GPU is required --- model inference does not run on CPU in
any practical configuration. Table~\ref{tab:hw} lists the rough VRAM and disk
budgets.

\begin{table}[H]
\centering\small
\renewcommand{\arraystretch}{1.25}
\begin{tabular}{>{\raggedright\arraybackslash}p{3.0cm}
                >{\raggedright\arraybackslash}p{6.2cm}
                >{\raggedright\arraybackslash}p{4.4cm}}
\toprule
\textbf{Mode} & \textbf{GPU / VRAM} & \textbf{Disk} \\
\midrule
Two-GPU &
  Card A $\ge 24$\,GB (embedder \code{gte-Qwen2-7B} fp16 ${\sim}15$\,GB $+$
  reranker ${\sim}2$\,GB $+$ transient MinerU); \newline
  Card B $\ge 24$\,GB (LLM \code{gemma4:31b} Q4 ${\sim}18$\,GB) &
  ${\sim}60$\,GB (weights $+$ images) \\
Single-GPU &
  One card $\ge 24$\,GB with a 12B LLM (\code{gemma3:12b} ${\sim}9$\,GB $+$
  small embedder $+$ reranker $+$ MinerU); \newline
  $\ge 40$\,GB if you run \code{gemma4:31b} on the same card &
  ${\sim}40$\,GB (weights $+$ images) \\
\bottomrule
\end{tabular}
\caption{Approximate hardware budgets. RAM $\ge 16$\,GB and a fast SSD are
recommended; the first start downloads several GB of model weights.}
\label{tab:hw}
\end{table}

\subsection{Software}
\begin{itemize}[leftmargin=1.4em,itemsep=2pt]
  \item \textbf{Git}.
  \item \textbf{Docker Engine} ($\ge 24$) with the \textbf{Compose v2} plugin
        (the \code{docker compose} subcommand), or Docker Desktop.
  \item \textbf{NVIDIA GPU driver} new enough for CUDA 12.x.
  \item \textbf{NVIDIA Container Toolkit} --- lets containers see the GPU.
\end{itemize}

\subsection{GPU + Docker setup by operating system}

\paragraph{Linux (Ubuntu/Debian) --- recommended.}
Install Docker Engine, then the NVIDIA Container Toolkit:
\begin{lstlisting}
# 1) Docker Engine + Compose plugin
curl -fsSL https://get.docker.com | sh
sudo usermod -aG docker "$USER"   # log out/in afterwards

# 2) NVIDIA Container Toolkit (official repo)
curl -fsSL https://nvidia.github.io/libnvidia-container/gpgkey \
  | sudo gpg --dearmor -o /usr/share/keyrings/nvidia-container-toolkit-keyring.gpg
curl -s -L https://nvidia.github.io/libnvidia-container/stable/deb/nvidia-container-toolkit.list \
  | sed 's#deb https://#deb [signed-by=/usr/share/keyrings/nvidia-container-toolkit-keyring.gpg] https://#g' \
  | sudo tee /etc/apt/sources.list.d/nvidia-container-toolkit.list
sudo apt-get update && sudo apt-get install -y nvidia-container-toolkit

# 3) Wire the toolkit into Docker and restart
sudo nvidia-ctk runtime configure --runtime=docker
sudo systemctl restart docker
\end{lstlisting}
Verify the GPU is visible inside a container:
\begin{lstlisting}
docker run --rm --gpus all nvidia/cuda:12.4.1-base-ubuntu22.04 nvidia-smi
\end{lstlisting}
You should see your GPU(s). (The NVIDIA \emph{driver} itself is installed on the
host the usual way for your distro, e.g.\ \code{sudo apt-get install nvidia-driver-550}.)

\paragraph{Windows 10/11 --- via WSL2.}
\begin{enumerate}[leftmargin=1.5em,itemsep=2pt]
  \item Install WSL2: \code{wsl --install} in an admin PowerShell, then reboot.
  \item Install the \textbf{NVIDIA Windows driver} that includes WSL2/CUDA
        support (the standard Game-Ready/Studio driver does). Do \emph{not}
        install a separate driver inside WSL.
  \item Install \textbf{Docker Desktop} and enable
        \emph{Settings $\rightarrow$ Resources $\rightarrow$ WSL integration}.
  \item In a WSL2 (Ubuntu) shell, verify:
        \code{docker run --rm --gpus all nvidia/cuda:12.4.1-base-ubuntu22.04 nvidia-smi}.
\end{enumerate}
Clone and run the project \emph{inside the WSL2 filesystem} (e.g.\ under
\code{\textasciitilde/} in Ubuntu, not \code{/mnt/c/...}) for correct file
permissions and speed.

\paragraph{macOS --- not supported for the GPU stack.}
Docker Desktop on macOS cannot pass through an NVIDIA GPU (there is no CUDA on
Mac). The embedder, reranker, LLM, and MinerU therefore cannot run locally. Use
a Linux host or a remote GPU server for the service; a Mac can act as a
\emph{client} (open the dashboard in a browser, optionally over an SSH tunnel:
\code{ssh -L 3006:localhost:3006 user@gpu-host}).

% =====================================================================
\section{Get the code}
% =====================================================================
\begin{lstlisting}
git clone https://github.com/desmondc65/NTU-DAVID-RAG.git
cd NTU-DAVID-RAG
\end{lstlisting}
All deployment files live under \code{docker/}.

% =====================================================================
\section{Configure secrets and ports}\label{sec:config}
% =====================================================================
The stack reads its configuration from \code{docker/.env}. Create it from the
template and set the two required secrets:
\begin{lstlisting}
cd docker
cp .env.example .env

# Generate strong secrets and paste them into .env:
openssl rand -hex 24    # -> AUTH_DB_PASSWORD
openssl rand -hex 32    # -> SESSION_SECRET
\end{lstlisting}

\noindent Key settings in \code{docker/.env} (defaults shown):
\begin{table}[H]
\centering\small
\renewcommand{\arraystretch}{1.2}
\begin{tabular}{>{\raggedright\arraybackslash}p{4.2cm}
                >{\raggedright\arraybackslash}p{3.0cm}
                >{\raggedright\arraybackslash}p{6.0cm}}
\toprule
\textbf{Variable} & \textbf{Default} & \textbf{Meaning} \\
\midrule
\code{UNIFIED\_PORT}        & \code{3006} & External dashboard port. \\
\code{EMBEDDING\_GPU}       & \code{1}    & Physical GPU index for the embedder/reranker. \\
\code{RAG\_WEB\_GPU}        & \code{0}    & GPU index for MinerU during ingestion. \\
\code{MINERU\_HYBRID\_BATCH\_RATIO} & \code{8} & Caps MinerU's batch size to avoid OOM on a shared card. \\
\code{EMBEDDING\_MODEL}     & \code{gte-Qwen2-7B} & Embedder (must match what populated the DB). \\
\code{AUTH\_DB\_PASSWORD}   & --- & \textbf{Required.} PostgreSQL password. \\
\code{SESSION\_SECRET}      & --- & \textbf{Required.} Session/CSRF signing key. \\
\code{ALLOW\_REGISTRATION}  & \code{true} & Set \code{false} after creating accounts. \\
\code{SESSION\_COOKIE\_SECURE} & \code{false} & Set \code{true} when behind HTTPS. \\
\bottomrule
\end{tabular}
\caption{Core environment variables (full list in \code{docker/.env.example}).}
\label{tab:env}
\end{table}

% =====================================================================
\section{Option A --- Two-GPU deployment (default)}\label{sec:twogpu}
% =====================================================================
The embedder/reranker run on \code{EMBEDDING\_GPU}; the answer LLM
(\code{gemma4:31b}) runs as a separate Ollama service on its own GPU.

\paragraph{Step 1 --- Start the answer LLM.}
The standalone Ollama compose pins the model to GPU index \code{2}; edit
\code{NVIDIA\_VISIBLE\_DEVICES}/\code{CUDA\_VISIBLE\_DEVICES} in the file to pick
another card.
\begin{lstlisting}
cd docker/local_llm
docker compose -f docker-compose.gemma4_31b_ollama.yml up -d
# First start pulls gemma4:31b (~19 GB) and warms it. API on :11434.
docker logs -f gemma4_ollama        # wait for "Model ready."
\end{lstlisting}

\paragraph{Step 2 --- Start the RAG stack.}
From \code{docker/}, with your \code{.env} in place:
\begin{lstlisting}
cd ..                 # back to docker/
./up.sh --build       # builds images, starts auth-db, embedding-server,
                      # rag-web, and the proxy (foreground; logs to log.txt)
# Detached instead:   ./up.sh -d --build
\end{lstlisting}
The first boot downloads the 7B embedder (${\sim}29$\,GB) and reranker into
\code{docker/../hf\_cache}; the \code{embedding-server} healthcheck has a
${\sim}10$\,minute grace window while this happens. \code{rag-web} starts only
once the embedder is healthy.

\paragraph{Step 3 --- Confirm the LLM connection.}
\code{rag-web} reaches the LLM via \code{docker/RAG/.env}
(\code{LOCAL\_LLM\_BASE\_URL}, default \code{http://host.docker.internal:11434/v1},
and \code{LLM\_MODEL\_NAME=gemma4:31b}). Adjust if your LLM runs elsewhere.

Continue at Section~\ref{sec:firstrun}.

% =====================================================================
\section{Option B --- Single-GPU deployment}\label{sec:onegpu}
% =====================================================================
A compose override (\code{docker/docker-compose.one-gpu.yml}) runs
\emph{everything} on one card: a lighter embedder, the reranker, MinerU, and an
in-stack Ollama LLM. Your default two-GPU setup is untouched --- this is an
opt-in launcher. Full notes: \code{docker/README.one-gpu.md}.

\paragraph{Step 1 --- (Optional) pick models and token.}
Defaults: embedder \code{embeddinggemma-300m}, LLM \code{gemma3:12b-it-qat}.
\code{embeddinggemma} is \emph{gated} on Hugging Face --- either accept its
license and set a token, or switch to a non-gated embedder. Put any overrides in
\code{docker/.env}:
\begin{lstlisting}
# docker/.env  (single-GPU knobs; all optional)
HF_TOKEN=hf_xxx                              # needed only for embeddinggemma
ONE_GPU_EMBEDDING_MODEL=BAAI/bge-large-en-v1.5   # non-gated alternative
ONE_GPU_LLM_MODEL=gemma3:12b-it-qat
# OLLAMA_CONTEXT_LENGTH=32768                 # cap the LLM context to fit VRAM
\end{lstlisting}

\paragraph{Step 2 --- Launch on a chosen card.}
\begin{lstlisting}
cd docker
GPU=0 ./up.one-gpu.sh --build     # GPU=<physical index> selects the card
# Detached:  GPU=0 ./up.one-gpu.sh -d --build
\end{lstlisting}
First boot builds the embedding image, pulls the LLM (${\sim}8$--$19$\,GB), and
loads the models; the in-stack \code{ollama} healthcheck waits for the pull to
finish before \code{rag-web} starts.

\paragraph{Note --- switching embedders is non-destructive.}
Each embedder writes to its own Qdrant collections
(\code{rag\_embeddings\_\_<model>}), so different-dimension models coexist; you
do not need to wipe the database when changing \code{ONE\_GPU\_EMBEDDING\_MODEL}.
The new model's collections start empty, so re-ingest your papers under it.

% =====================================================================
\section{First run and verification}\label{sec:firstrun}
% =====================================================================
\begin{enumerate}[leftmargin=1.5em,itemsep=3pt]
  \item \textbf{Watch startup.} \code{docker compose ps} shows service health;
        \code{docker logs -f embedding-server} shows the embedder loading.
        Expect several minutes on a cold cache.
  \item \textbf{Open the dashboard:} \code{http://<host>:3006/}. Over SSH:
        \code{ssh -L 3006:localhost:3006 user@host} then browse
        \code{http://localhost:3006/}.
  \item \textbf{Create an account.} Register on the login page (open
        registration is on by default). Set \code{ALLOW\_REGISTRATION=false} in
        \code{.env} and restart once your accounts exist.
  \item \textbf{Ingest a paper.} Upload a PDF (and, optionally, its Fortran
        source). Ingestion is synchronous --- PDF parsing, enrichment, code
        digestion, chunking, embedding --- and can take several minutes per
        paper on the first run.
  \item \textbf{Ask a question.} Each paper's vectors are isolated per user;
        answers cite paper titles and are grounded in retrieved context.
\end{enumerate}

% =====================================================================
\section{Day-to-day operations}
% =====================================================================
Run these from \code{docker/} (add \code{-f docker-compose.yml -f
docker-compose.one-gpu.yml} for single-GPU mode, or just use the matching
\code{up} script).

\begin{table}[H]
\centering\small
\renewcommand{\arraystretch}{1.25}
\begin{tabular}{>{\raggedright\arraybackslash}p{7.2cm}
                >{\raggedright\arraybackslash}p{6.6cm}}
\toprule
\textbf{Task} & \textbf{Command} \\
\midrule
Start (two-GPU)        & \code{./up.sh -d} \\
Start (single-GPU)     & \code{GPU=0 ./up.one-gpu.sh -d} \\
Status / health        & \code{docker compose ps} \\
Follow a service log   & \code{docker logs -f rag-web} \\
Rebuild after a code change & \code{./up.sh -d --build} \\
Stop (keep data)       & \code{docker compose down} \\
Stop single-GPU stack  & \code{docker compose -f docker-compose.yml -f docker-compose.one-gpu.yml down} \\
GPU usage              & \code{nvidia-smi} \\
\bottomrule
\end{tabular}
\caption{Common operational commands.}
\label{tab:ops}
\end{table}

\noindent Persistent state lives in the repository tree and survives
\code{down}: \code{db/} (Qdrant vectors), \code{data/} (uploaded papers),
\code{auth\_db/} (PostgreSQL), and \code{hf\_cache/} (downloaded model weights).

% =====================================================================
\section{Troubleshooting}
% =====================================================================
\begin{itemize}[leftmargin=1.4em,itemsep=3pt]
  \item \textbf{``could not select device driver \dots gpu''} --- the NVIDIA
        Container Toolkit is not configured. Re-run
        \code{sudo nvidia-ctk runtime configure --runtime=docker} and restart
        Docker; confirm with the \code{--gpus all} test in Section~\ref{sec:req}.
  \item \textbf{Port 3006 already in use} --- set a different
        \code{UNIFIED\_PORT} in \code{.env} and restart.
  \item \textbf{\code{embedding-server} stuck ``starting''} --- it is still
        downloading the 7B weights (${\sim}29$\,GB). Watch
        \code{docker logs -f embedding-server}; the healthcheck allows
        ${\sim}10$\,min.
  \item \textbf{GPU out-of-memory during ingestion} --- lower
        \code{MINERU\_HYBRID\_BATCH\_RATIO} (8 $\rightarrow$ 4 $\rightarrow$ 2),
        and in single-GPU mode lower \code{OLLAMA\_CONTEXT\_LENGTH}.
  \item \textbf{Ingestion times out (504)} --- very large Fortran files make a
        single upload exceed the proxy timeout; the bundled Nginx/Gunicorn
        timeouts are 3600\,s, which suits large files but still bounds them.
  \item \textbf{Empty answers / LLM errors} --- verify the LLM is reachable:
        \code{curl http://localhost:11434/api/tags} should list the model named
        in \code{LLM\_MODEL\_NAME}.
  \item \textbf{Gated-model 401 (single-GPU)} --- set \code{HF\_TOKEN} in
        \code{.env} for \code{embeddinggemma}, or use the non-gated
        \code{BAAI/bge-large-en-v1.5}.
\end{itemize}

% =====================================================================
\section{Quick reference}
% =====================================================================
\begin{lstlisting}
# ---- Two-GPU (default) -------------------------------------------------
git clone https://github.com/desmondc65/NTU-DAVID-RAG.git
cd NTU-DAVID-RAG/docker && cp .env.example .env      # set AUTH_DB_PASSWORD + SESSION_SECRET
( cd local_llm && docker compose -f docker-compose.gemma4_31b_ollama.yml up -d )
./up.sh -d --build
# open http://<host>:3006/

# ---- Single-GPU -------------------------------------------------------
cd NTU-DAVID-RAG/docker && cp .env.example .env      # set secrets (+ HF_TOKEN or bge-large)
GPU=0 ./up.one-gpu.sh -d --build
# open http://<host>:3006/
\end{lstlisting}

\end{document}
